% ======================================================================
% Chapter 22: Quantum Field Theory at the Frontier
% (book pages 781-823; PDF pages 804-823) -- The Epilogue
% ======================================================================
\chapter{Quantum Field Theory at the Frontier}

In this textbook we have surveyed the most important ideas of quantum field
theory.  Working from the basic concepts that come from fusing relativity,
quantum mechanics, and fields, we have built an elaborate structure, which
includes such remarkable elements as coupling constant renormalization and
non-Abelian gauge fields.  We have seen at many points that the strange and
abstract elements of this structure actually intersect with observation and
even produce explanations for unexpected aspects of the behavior of
elementary particles.

In the course of our study, we have arrived at a complete theory of the
strong, weak, and electromagnetic interactions of elementary particles.  Each
element of this theory has been described as a quantum field theory, and
these quantum field theories have turned out to have very similar structure
as gauge theories coupled to fermions.  At various points in our discussion,
we have noted that these theories have passed stringent quantitative
experimental tests.  We have not had space to describe the wide variety of
experiments that contribute to our faith in these theories, but to day almost
all particle physicists consider this $SU(3)\times SU(2)\times U(1)$ gauge
theory as established.  In fact, most of these people refer to this theory
condescendingly as ``the standard model''.

Though the best data to support the standard model have come from experiments
of the past five years, the ideas behind it are much older.  Most of the
theoretical developments described in this book were concluded in the 1970s,
a generation removed from the current frontier of physics.  But this does not
mean that quantum field theory is irrelevant to that frontier, any more than
quantum mechanics and electrodynamics have lost their relevance after many
years of exploration.  On the contrary, the theory of elementary particles---like
other areas of physics that depend on quantum fluctuations in continua---still
holds deep mysteries, and quantum field theory remains the principal tool for
exploration of these questions.

In this final chapter, we will flash forward to the present day and discuss
the relevance of quantum field theory to current questions in the physics of
the fundamental interactions.  We will present what are, in our view, the
outstanding unsolved problems of elementary particle physics and describe how
quantum field theory is being used to confront these problems.  Many of these
applications involve aspects of quantum field theory that are beyond the
scope of this book.  These include the use of quantum field theory in the
regime beyond the reach of perturbation theory and the use of quantum field
theory to explore the special properties of theories with higher spin and
local symmetry.  In these areas our discussion will be mainly qualitative, but
we will give references that provide points of entry into each of these
subjects.

It should be obvious that our discussion in this final chapter will express
our personal opinions and by no means represents the consensus of experts in
quantum field theory.  In addition, any collection of ``current problems'' in
a rapidly developing area of research should quickly become dated.  In fact,
we hope the readers of this book will quickly make this chapter obsolete by
solving the problems that we highlight here.

\section{Strong Strong Interactions}

One paradoxical aspect of our discussion of the strong interactions is that
all of our concrete results were obtained by assuming that these
interactions are weak.  At large momentum transfer, we argued, this
assumption is actually valid due to asymptotic freedom.  Still, it is
uncomfortable that we have left the most obvious questions about strongly
interacting particles---for example, their masses and low-energy
interactions---in a mysterious regime excluded from our analysis.

To work with QCD in the region where the strong interactions are strong, we
need to answer three questions: First, how do we describe the forces that
bind quarks together into hadrons?  Second, what is an appropriate
description of the quark-antiquark and three-quark systems bound by those
forces?  And finally, how do we compute scattering amplitudes and matrix
elements of currents using these bound states?

At this moment, there is no derivation of the complete force between quarks
from the QCD Lagrangian.  Explicit calculations can be done only in the
limits of weak and strong coupling.  In the weak-coupling limit, the result
is a Coulomb potential with an asymptotically free coupling constant.  The
strong coupling limit, on the other hand, gives a linear potential which
confines color in the way that we described, but did not derive, at the end
of Section~17.1.  The derivation of this result is quite unusual and brings
in a new set of mathematical methods.

So far in this book, we have not discussed a strong coupling approximation to
a quantum field theory.  There is a simple reason for this: In a quantum
field theory in which the coupling $g^2$ is very large, the elementary
particles or their bound states typically acquire masses that grow with
$g^2$.  For $g^2\to\infty$, these masses become comparable to the cutoff
$\Lambda$ and the field theory ceases to have a local continuum
description.

Wilson proposed to solve this problem in a radical way, by replacing
spacetime with a lattice of discretely spaced points.  Such a lattice is
easiest to visualize in Euclidean spacetime, and so we can use a functional
integral over fields on a lattice to approximate Euclidean Green's functions.
Such a theory can have a well-defined strong coupling limit.  A theory of
this type is very similar to a lattice model of a magnetic system.

In fact, we can understand this construction of a quantum field theory by
using the concepts of Chapter~13.  A lattice theory with fluctuating spin
variables at each lattice site is described in the large by a quantum field
theory of scalar fields with the symmetry of the underlying spin variables.
Typically, the strong-coupling limit of the quantum field theory corresponds
to the high-temperature limit of the magnet, in which the correlation length
is much smaller than the lattice spacing.  Decreasing the coupling constant
corresponds to decreasing the temperature.  Eventually, the coupling constant
comes close to a fixed point of the renormalization group, and one can use
this fixed point to define a limit of the lattice functional integral in
which the lattice spacing is taken to zero.

To build a lattice model of the strong interactions, one needs to find a set
of variables on the discrete lattice that correspond in the large to
non-Abelian gauge fields.  Wilson proposed that the fundamental variables for
such a theory should be the line elements from one lattice vertex $v_1$ to a
neighboring vertex $v_2$,
\begin{equation}
  U(v_1, v_2) ,
  \label{eq:22.1}
\end{equation}
where each $U$ is an element of the gauge group $G$.  To construct the
lattice gauge theory with gauge group $G$, one should integrate over a finite
group transformation $U$ for each link of the lattice.  Taking a product of
these $U$ matrices around a closed path, one can construct gauge-invariant
observables, just as we did in Section~15.3.  An appropriate Lagrangian can
also be constructed as a sum of gauge-invariant products of the $U$ matrices
about elementary closed loops of the lattice.%
\footnote{This construction was introduced by K.~Wilson, \emph{Phys.\ Rev.\ D}
\textbf{10}, 2445 (1974).  The construction is described pedagogically in
M.~Creutz, \emph{Quarks, Gluons, and Lattices} (Cambridge University Press,
Cambridge, 1983).}

In a spin system, the defining property of the high-temperature phase is the
exponential decay of correlations
\begin{equation}
  \langle s(0)\cdot s(x)\rangle \sim \exp[-|x|/\xi]
  \label{eq:22.2}
\end{equation}
as $|x|\to\infty$.  The analogous property of the gauge-invariant correlation
function of $U$ matrices around a closed path $P$ is
\begin{equation}
  \bigl\langle \mathrm{tr}\prod_P U \bigr\rangle \sim e^{-\sigma A},
  \label{eq:22.3}
\end{equation}
where $A$ is the area spanned by the path.  This behavior is in fact seen
explicitly in the expansion of Wilson's lattice gauge theory for strong
coupling.  A pair of color sources that stand a distance $R$ apart for a
Euclidean time $T$ are represented by a large rectangular loop of width $R$
and length $T$.  For such a loop, we can compare the result \eqref{eq:22.3}
to the expression for the energy of an excited state in Euclidean time,
\begin{equation}
  \langle \exp[-H_E T] \rangle \sim \exp[-E\,T] .
  \label{eq:22.4}
\end{equation}
Then we see that static sources of gauge charge, in the strong-coupling
limit, are attracted to one another by a potential energy
\begin{equation}
  V(R) \sim \sigma R
  \label{eq:22.5}
\end{equation}
at sufficiently large $R$, where $\sigma$ is the string tension.  Similarly,
when one introduces dynamical quarks into a lattice gauge theory and studies
their properties in the strong-coupling limit, configurations with large
separation of color sources are suppressed in the Euclidean functional
integral by factors of the form of \eqref{eq:22.3}.  The strong coupling
limit then predicts the permanent confinement of quarks into color-singlet
bound states.

The argument we have just given applies equally well to gauge theories based
on Abelian or non-Abelian symmetry groups.  But non-Abelian gauge theories
have the important additional property of asymptotic freedom.  In this
context, that implies that a theory with weak coupling at short distances
flows to a theory with strong coupling at large distances.  If we imagine
integrating out short-distance degrees of freedom as we described in
Section~12.1, and if there is no zero of the $\beta$ function or other
barrier to the renormalization group flow, we should eventually arrive at an
effective theory for which the strong-coupling expansion is a good
approximation.  Thus, in the particular case of non-Abelian gauge theories,
asymptotic freedom allows us to connect a short-distance picture based on
free quarks and gluons to a large-distance picture based on color
confinement.

It would be wonderful if the strong-coupling picture that we have described
led to mathematical equations in continuum spacetime describing the motion
of permanently confined quarks and antiquarks.  Many authors have tried to
write such equations by imagining the area suppression of the Wilson loop
correlation function \eqref{eq:22.3} to result from a physical surface that
spans the loop.  For the large rectangular loop associated with color
sources, this surface can be interpreted physically as the lines of color
electric flux that run from one source to the other (as in
Fig.~\ref{fig:17.1}), swept out through Euclidean time.  At one moment of
Euclidean time, this surface can be idealized as an abstract one-dimensional
excitation, often called a \emph{string}.  Unfortunately, the quantum
properties of an idealized string turn out to be very complicated, since each
small element of the string must be considered as an independent quantum
degree of freedom.  The only systems of string equations that have actually
been solved have bizarre features, including unwanted massless particles.  Up
to now, no one has succeeded in writing an equation for the quark-confining
string that is useful for quantitative calculations of quark bound
states.%
\footnote{For one approach to color confinement from a picture involving
Wilson loops and strings, see A.~A.~Migdal, \emph{Phys.\ Repts.} \textbf{102},
199 (1983).}

However, the lattice regularization of a non-Abelian gauge theory suggests
another approach to quantitative calculations in strong-interaction theory.
By approximating QCD by a lattice gauge theory with a nonzero lattice
spacing and a finite spacetime volume, we reduce the functional integral to a
finite number of bounded integrations, that is, an integral over $SU(3)$
group matrices for each of the finite number of links in the lattice.  A
lattice of size, for example, $20^4$ allows the lattice spacing to be smaller
than the size of a hadron while the full size of the lattice is much larger
than a hadronic radius.  Then one can compute correlation functions by
evaluating the integrals numerically, by the Monte Carlo method.  Since the
functional integral with a finite lattice spacing is related to the original
functional integral with zero lattice spacing by integrating out short-distance
degrees of freedom, the lattice approximation can be systematically improved
by computing the short-distance effects perturbatively, using asymptotic
freedom to justify a weak-coupling analysis.%
\footnote{For an introduction to numerical lattice gauge theory, see
\emph{From Actions to Answers}, T.~DeGrand and D.~Toussaint, eds.\ (World
Scientific, Singapore, 1990).}

This numerical method has now become the principal theoretical tool for
quantitative calculations in hadron physics.  This method currently gives the
masses of the low-lying mesons and baryons to accuracies of 10--20\%; it also
allows the calculation of weak interaction matrix elements of hadrons at the
25\% level.  As computers become more powerful, this numerical method can be
pushed to higher accuracy.

Eventually, it will be interesting to ask whether these nonperturbative
numerical calculations are consistent with our precision knowledge of the
perturbative region of QCD.  At the time of this writing, the first such
comparison has been made: We have listed in Table~\ref{tab:17.1} a value of
$\alpha_s$ from $J/\psi$ and $\Upsilon$ spectroscopy.  In this calculation,
the experimentally determined masses of $c\bar{c}$ and $b\bar{b}$ bound states
are compared to values computed numerically with lattice regularization.  The
comparison of these values gives the required bare coupling constant of the
lattice theory, which can be converted to a value of $\alpha_s(m_Z)$ in the
convention of the table.  The resulting estimate for $\alpha_s(m_Z)$ does
agree reasonably well with purely perturbative determinations.

What is the future of nonperturbative calculations in hadron physics?  On the
one hand, we expect to see further development of numerical lattice methods.
These methods have hardly begun to address problems of hadron-hadron
scattering and multiparticle matrix elements, and this seems an important
direction for the future.  In addition, these methods should eventually
supply an engineering understanding of hadrons at the percent level or
better.  On the other hand, we hope also to see a quantitative continuum
approach to hadron structure, in which dynamical quarks interact with some
appropriate type of string degrees of freedom.

\section{Grand Unification and its Paradoxes}

If we put aside our questions about the low-energy, nonperturbative behavior
of QCD, the $SU(3)\times SU(2)\times U(1)$ gauge theory gives an apparently
complete description of elementary particle interactions at those energies
that we have probed experimentally.  But what happens beyond our current
reach?  Does this theory need modification, or could it continue to be valid
at much higher energies?

The $SU(3)\times SU(2)\times U(1)$ gauge theory contains three independent
gauge coupling constants, and the observed values of these couplings are
larger for the larger components of the gauge group.  This pattern can be
explained by a bold hypothesis about the behavior of the gauge couplings at
very high energy.  If at some very large energy scale, these three couplings
were equal, the values of the $SU(3)$ and $SU(2)$ couplings would increase at
smaller momentum scales due to their asymptotically free renormalization
group equations, while the value of the $U(1)$ coupling would decrease,
resulting in the observed pattern of couplings at low energies.

An even bolder hypothesis would be that the three gauge symmetries are
subgroups of a single large symmetry group, which is spontaneously broken at
very high energy scales.  The simplest choice for this larger symmetry is
$SU(5)$.  In that theory, the coupling constants of $SU(3)\times SU(2)\times
U(1)$ have the following relation to the underlying $SU(5)$ coupling at the
scale of $SU(5)$ breaking:
\begin{equation}
  g_3 = g_2 = \sqrt{\frac{5}{3}}\,g_1 .
  \label{eq:22.6}
\end{equation}
The idea that the $SU(3)\times SU(2)\times U(1)$ gauge group is embedded in a
larger simple group is known as \emph{grand unification}; the particular
choice of $SU(5)$ as the unifying group is due to Georgi and Glashow.%
\footnote{The theory of grand unification is reviewed in P.~Langacker,
\emph{Phys.\ Repts.} \textbf{72}, 185 (1981), and in P.~Langacker,
\emph{Phys.\ Repts.} \textbf{109}, 145 (1984).}  The observed quarks and
leptons can be seen to fit neatly into an anomaly-free chiral representation
of $SU(5)$; this embedding leads to a natural explanation of the fractional
charges of quarks.%

Within this framework, we can extrapolate the values of the three coupling
constants from the energy scale of $m_Z$ upward.  The result of this
extrapolation is shown as the solid lines in Fig.~22.1.  The
coupling constants do come close together at very high energies, though they
do not actually meet.  The dashed lines in the figure show the evolution with
a modified set of renormalization group equations, to be explained in
Section~22.4; with this choice, the three couplings meet accurately within
their current uncertainties.  In any event, the evolution of coupling
constants occurs on a logarithmic scale in energy, so the energy at which the
couplings would meet is determined only logarithmically.  The result of the
extrapolation of the three couplings using the standard model is a grand
unification scale of order
\begin{equation}
  M_U \sim 10^{15}\ \mathrm{GeV} .
\end{equation}

Grand unification makes a number of dramatic predictions.  Since the quarks
and leptons of each generation are placed in a single representation of
$SU(5)$, the theory predicts transitions between quarks and leptons.  These
are mediated by new gauge bosons, the $X$ and $Y$ bosons, which are
extremely massive, of order $M_U$.  These interactions violate baryon
number, leading to proton decay.  The predicted lifetime of the proton is
\begin{equation}
  \tau_p \sim \frac{M_U^4}{\alpha_U^2 m_p^5} \sim 10^{30}\ \mathrm{years} ,
  \label{eq:22.7}
\end{equation}
which is many orders of magnitude larger than the age of the universe.  The
absence of observed proton decay places stringent bounds on the grand
unification scale and rules out the simplest $SU(5)$ model.  Nevertheless,
grand unification remains an important and attractive idea.

\subsection{The Gauge Hierarchy Problem}

Grand unification leads to a serious conceptual problem concerning the Higgs
boson mass.  In order to produce a vacuum expectation value of the right
size to give the observed $W$ and $Z$ boson masses, the Higgs scalar field
must obtain a negative mass term, of the size
\begin{equation}
  -\mu^2 \sim -(100\ \mathrm{GeV})^2 .
  \label{eq:22.8}
\end{equation}
Unfortunately, the $(\mathrm{mass})^2$ of a scalar field receives additive
renormalizations.  In a theory with cutoff scale $\Lambda$, $\mu^2$ can be
much smaller than $\Lambda^2$ only if the bare mass of the scalar field is of
order $-\Lambda^2$, and this value is canceled down to $-\mu^2$ by radiative
corrections.  If we envision that our theory of Nature contains the very
large scales of grand unification, we must take seriously the idea that the
appropriate value to take for $\Lambda$ in this discussion is $10^{16}$ GeV
or larger.  This seems to require dramatic and even bizarre cancellations in
the renormalized value of $\mu^2$.

We met a situation of this type in the theory of phase transitions.  At zero
temperature, a ferromagnet typically has a spin expectation value of the
order of the underlying atomic parameters.  As the temperature is raised, or
as some other variable in the system is changed, the magnetization
decreases.  Finally, by fine adjustment of the temperature, we can arrive at
a situation where the system approaches a critical point.  In the very near
vicinity of this point, the expectation value of the spin field is much
smaller than the value predicted from atomic parameters, and the system is
described by an approximately massless continuum scalar field theory.

In statistical mechanics, this picture of the light scalar field makes sense
because there is an experimenter sensitively adjusting a dial.  In the theory
of weak interactions, there is no one obviously making a fine adjustment that
gives the $(\mathrm{mass})^2$ of the Higgs boson a value 28 orders of
magnitude or more below its natural value.  Thus, it is a mystery why the
Higgs boson mass should be so small compared to the grand unification scale.
Particle physicists refer to this question as the \emph{gauge hierarchy
problem}.

How can one naturally arrange a Higgs field mass term to be so much smaller
than the underlying mass scale of the fundamental interactions?  One
possible strategy would be to arrange for a symmetry of the fundamental
Lagrangian that forbids the Higgs boson mass term and that is very weakly
broken.  This idea turns out to be very difficult to implement.  To build a
theory of this type, one would need to create a scalar field theory in which
additive radiative corrections to the Higgs boson mass must cancel to any
foreseeable order in perturbation theory.  But the Higgs mass term is very
simple in form,
\begin{equation}
  \Delta\mathcal{L} = \mu^2|\phi|^2 ,
  \label{eq:22.9}
\end{equation}
and it is hard to imagine any principle that would keep this term from being
generated by radiative corrections.  There is one proposal for a symmetry
with this property, but it requires the introduction of a profound principle
called \emph{supersymmetry} that entails deep modifications of fundamental
physics.  In particular, it requires a large number of new elementary
particles, some of which should have masses below 1000 GeV, within the reach
of the next generation of accelerators.  We will discuss this possibility
further in Section~22.4.

In this discussion, the problem of the Higgs mass stemmed from the hypothesis
that the Higgs boson was an elementary particle.  An alternative viewpoint,
already suggested at the end of Section~20.2, is that the Higgs boson is a
composite state bound by a new set of interactions.  This idea also leads to
observable experimental consequences, since the mass scale of these new
interactions must be close to the weak interaction mass scale.  In the
simplest theories of this type, the symmetry breaking of the Higgs sector is
modeled on the dynamical chiral symmetry breaking of the strong
interactions, which we discussed in Section~19.3.  The new strong
interactions required by the theory lead to a spectrum of new particles with
masses of about 1000 GeV.%
\footnote{The properties of these models of the Higgs sector, known to
specialists as technicolor models, are described in R.~Kaul,
\emph{Rev.\ Mod.\ Phys.} \textbf{55}, 449 (1983) and K.~D.~Lane, in
\emph{The Building Blocks of Creation}, S.~Raby, ed.\ (World Scientific,
1993).}  Thus, the two conflicting hypotheses on the nature of the sector
that breaks $SU(2)\times U(1)$ both lead to new phenomena observable at
future accelerators, and possibly even at present ones.

Just as these two different theories of the Higgs sector present completely
different answers to the question of why the weak-interaction symmetry
$SU(2)\times U(1)$ should be spontaneously broken, they also imply completely
different answers to the question of the origin of the quark and lepton
masses.  In a model in which the Higgs field is elementary, the quark and
lepton masses are derived from the renormalizable couplings of fermions to
the Higgs field.  These couplings would presumably be part of the grand
unification and could be predicted only by theories that made explicit
reference to the grand unification scale.  In principle, the knowledge of
these couplings could give us clues as to the details of the grand
unification.  Even if the Higgs field is composite, we cannot avoid the fact
that the generation of masses for the quarks and leptons requires the
breaking of $SU(2)\times U(1)$.  Thus, these mass terms must arise from
couplings of the quarks and leptons to the Higgs sector of interactions.  In
this class of models, the interactions leading to the quark and lepton masses
must arise at energies close to the scale of the Higgs sector strong
interaction and may eventually be observable experimentally.

From either viewpoint, it is still mysterious why the spectrum of quarks and
leptons covers 5 orders of magnitude, from the electron at 0.5 MeV to the
top quark at 175 GeV.  It is also not understood what gives rise to the
pattern of quark mixings encoded in the CKM matrix and the magnitude of
$CP$ violation.  Even with detailed confirmation of the standard model,
these questions seem today very far from solution.

\subsection{The Cosmological Constant Problem}

The enormous mass scale of grand unification can also enter one more physical
quantity, one that poses an even greater paradox than that of the Higgs
boson mass.  When we first quantized a field in Section~2.3, we discovered
that the energy density of the vacuum in free scalar field theory received an
infinite positive contribution from the zero-point energies of the various
modes of oscillation.  With a cutoff scale $\Lambda$, this zero-point energy
is given roughly by
\begin{equation}
  \langle 0|H|0\rangle \sim \Lambda^4 .
  \label{eq:22.10}
\end{equation}
At many other points in our discussion, we found similarly large
contributions to the vacuum energy.  The filling of the Dirac sea in the
quantization of the free fermion theory led to a downward shift in the vacuum
energy with a similar ultraviolet divergence.  Spontaneous symmetry breaking
gives a finite but still possibly large shift in the vacuum energy density,
\begin{equation}
  \Delta\langle 0|H|0\rangle \sim -cv^4 ,
  \label{eq:22.11}
\end{equation}
with dimensionless $c$, for a field vacuum expectation value $v$.  The
spontaneous breaking of the weak interaction $SU(2)\times U(1)$ symmetry and
of the strong interaction chiral symmetry both would be expected to shift the
vacuum energy density in this way.

In elementary particle physics experiments, this shift of the vacuum energy
is unobservable.  Experimentally measured particle masses, for example, are
energy differences between the vacuum and certain excited states of $H$, and
the absolute vacuum energy cancels out in the calculation of these
differences.  However, there is a way that the absolute vacuum energy could
potentially be observed, through the coupling of the vacuum energy to
gravity.  According to Einstein, the energy-momentum tensor of matter
$\Theta^{\mu\nu}$ is the source of the gravitational field.  A vacuum energy
density $\langle 0|H|0\rangle = \Lambda$ contributes to this source a term
\begin{equation}
  \Theta^{\mu\nu} \supset \Lambda\,g^{\mu\nu} ,
  \label{eq:22.12}
\end{equation}
where the first term on the right is subtracted to have zero vacuum
expectation value.  The vacuum energy term has the form of Einstein's
\emph{cosmological constant} and thus potentially affects the expansion of
the universe.

In fact, measurements of the cosmological expansion exclude a large
cosmological constant.  The current limit is
\begin{equation}
  \Lambda < 10^{-29}\ \mathrm{g/cm}^3 \sim (10^{-11}\ \mathrm{GeV})^4 .
  \label{eq:22.13}
\end{equation}
We have no understanding of why $\Lambda$ is so much smaller than the vacuum
energy shifts generated in the known phase transitions of particle physics,
and so much again smaller than the underlying field zero-point energies.  The
discrepancy in $\Lambda$ between the experimental bound \eqref{eq:22.13} and
naive intuition is 120 orders of magnitude!  The solution to this problem may
come from one of many sources.  It may be that the formalism of the quantum
field theory of gravity requires that the vacuum energy be subtracted from
the energy-momentum tensor that appears in Einstein's equations of gravity.
It may be that there is a new physical mechanism coming from particle physics
or from gravity itself that sets the total vacuum energy to zero.  Or it may
be that the overall scale of energy-momentum is genuinely ambiguous and is
set by a cosmological boundary condition.  At this moment, all of these
possibilities are just guesses.  All we know for certain is that the
unification of quantum field theory and gravity cannot be straightforward,
that there is some important concept still missing from our current
understanding.%
\footnote{The cosmological constant problem and a variety of unsuccessful
solutions are reviewed in S.~Weinberg, \emph{Rev.\ Mod.\ Phys.} \textbf{61},
1 (1989).}

\section{Exact Solutions in Quantum Field Theory}

From the idea of grand unification, with its great promise and mystery, we
turn to the study of model quantum field theories that are so simple that
they can be solved exactly.  Throughout this book, we have stressed the
intrinsic complexity of quantum field theory and the importance of using
perturbation theory as a replacement for exact knowledge.  But there are a
variety of quantum field theories for which exact solutions are known.  In
this section, we will describe some of these and review the insights we have
gained from them.

In searching for exact solutions to quantum field theory models, there is no
reason to restrict our attention to four-dimensional spacetime.  In fact, we
have seen examples of two-dimensional theories with similar complexity of
renormalization and short-distance behavior.  At the same time, these
theories occupy a one-dimensional space, and their degrees of freedom can be
visualized as links in a chain.  This allows some powerful simplifications.

In our discussion of the axial anomaly in two dimensions in Section~19.1, we
showed that the photon of two-dimensional massless QED becomes a massive
boson.  More detailed examination of this theory shows that this boson is a
noninteracting particle.  The theory is originally formulated in terms of
fermions, interacting through Coulomb forces.  However, it is possible to
exactly rewrite the theory as a theory of a scalar field that creates and
destroys fermion-antifermion pairs.  Heuristically, a particle and an
antiparticle moving down the light-cone in one-dimensional space do not
separate and therefore comprise one bosonic degree of freedom.  In a wide
class of models, the bosonic theories rewritten in this way are free-field
theories.  A remarkable model of this type is the \emph{Thirring model},
\begin{equation}
  \mathcal{L} = \bar{\psi}i\slashed{\partial}\psi
    - \frac{1}{2}g(\bar{\psi}\gamma^\mu\psi)^2 ,
  \label{eq:22.14}
\end{equation}
in two dimensions.  In this model, the replacement of the fermion field by a
boson field leads to a free field theory.  Using this field theory, one can
compute correlation functions of fermion bilinears explicitly and show
directly that these operators have anomalous dimensions.  In
renormalization-group language, the model contains a line of fixed points
parametrized by the coupling constant $g$.%
\footnote{For an introduction to these models, see S.~Coleman,
\emph{Phys.\ Rev.\ D} \textbf{11}, 2088 (1975), \emph{Ann.\ Phys.} \textbf{101},
239 (1976).}

A more general class of two-dimensional models can be solved by visualizing
them in a Hamiltonian picture as a one-dimensional chain of coupled field
operators.  The prototype of this method is a problem in the statistical
mechanics of magnets, the one-dimensional chain of coupled spins.  Consider
a long chain of $N$ discrete sites, with a spin-$1/2$ system at each site.
The Pauli sigma matrices $\bm{\sigma}_i$ act on the two-dimensional Hilbert
space at the site $i$.  The Hamiltonian for the spin chain is then
\begin{equation}
  H = \sum_i \Bigl\{ -J\bm{\sigma}_i\cdot\bm{\sigma}_{i+1} \Bigr\} .
  \label{eq:22.15}
\end{equation}
Since
\begin{equation}
  \bm{\sigma}_i\cdot\bm{\sigma}_{i+1} = 2\bigl( \sigma^+_i\sigma^-_{i+1}
    + \sigma^-_i\sigma^+_{i+1} \bigr) + \sigma^z_i\sigma^z_{i+1} ,
  \label{eq:22.16}
\end{equation}
this Hamiltonian conserves the number of up spins.  The state with all spins
down is an eigenstate of the Hamiltonian, and the states with one spin up in
a state of definite momentum are also eigenstates.  In 1934, Bethe analyzed
the problem of two spins up and computed their $S$-matrix.  He then
discovered an amazing fact, that by multiplying the $S$-matrices for the
two-spin problem, he could find the exact eigenstates of the Hamiltonian for
any number of spins up.  By considering $N/2$ spins up, he found the ground
state of the system.  This technique, now known as \emph{Bethe's ansatz}, has
been used to solve a wide variety of one-dimensional problems in condensed
matter physics and quantum field theory.  For example, this technique has
been used by Andrei and Lowenstein to solve the Gross-Neveu model presented
in Problem~11.3 and to demonstrate that the spectrum of this model has the
properties expected from asymptotic freedom.%
\footnote{For an introduction to Bethe's ansatz and its generalizations, see
N.~Andrei, K.~Furuya, and J.~H.~Lowenstein, \emph{Rev.\ Mod.\ Phys.}
\textbf{55}, 331 (1983), L.~D.~Faddeev, in \emph{Recent Advances in Field
Theory and Statistical Mechanics}, J.~B.~Zuber and R.~Stora, eds.
(North-Holland, Amsterdam, 1984), and R.~J.~Baxter, \emph{Exactly Solved
Models in Statistical Mechanics} (Academic Press, London, 1982).}

Even where it is not possible to solve a model for all values of its
parameters, it is sometimes possible to find exact information about
two-dimensional models at points where they contain massless fields.  It is
well known that a variety of classical two-dimensional partial differential
equations can be solved by exploiting techniques of complex variables.  For
example, the two-dimensional Laplace equation $\nabla^2\phi = 0$ is
invariant to conformal mappings $z\to w(z)$, where $z = x + iy$.
Two-dimensional quantum field theories with massless particles often have
this \emph{conformal symmetry} at the classical level, though generically it
is anomalous.  In special systems, however, these anomalies vanish and the
quantum theory is invariant to conformal mapping.  These theories typically
contain operators with anomalous dimensions, indicating that each such theory
is a new, nontrivial fixed point of the renormalization group.  The conformal
symmetry of the theory can be used to compute these anomalous dimensions.

As an example of this class of theories, consider the two-dimensional
nonlinear sigma model in which the basic field is not a unit vector, as we
discussed in Section~13.3, but rather a unitary matrix of a Lie group $G$.
The Lagrangian of this theory is
\begin{equation}
  \mathcal{L} = \frac{1}{4\lambda^2}\ \mathrm{tr}\bigl[ \partial_\mu U^{-1}
    \partial^\mu U \bigr] ,
  \label{eq:22.17}
\end{equation}
where $U$ is a field taking values in the group $G$.  Like the theory of
Section~13.3, this model is asymptotically free.  However, Witten has shown
that, by adding to this Lagrangian a particular perturbation of a rather
complicated form first written by Wess and Zumino, one can find a fixed
point of the renormalization group with manifest $G\times G$ global
symmetry.  This theory is conformally invariant, and all operator correlation
functions can be computed using the conformal symmetry.%
\footnote{For an introduction to conformal field theory, see
A.~A.~Belavin, A.~M.~Polyakov, and A.~B.~Zamolodchikov,
\emph{Nucl.\ Phys.\ B} \textbf{241}, 333 (1984), and J.~L.~Cardy,
in \emph{Phase Transitions and Critical Phenomena}, C.~Domb and
J.~L.~Lebowitz, eds.\ (Academic Press, 1987).}

One result of the nonperturbative exploration of quantum field theory was the
realization that field theories can contain particle states that are not
simply related to the quanta of the original fields.  In the weak-coupling
limit of a quantum field theory, such new states can appear as new solutions
of the classical field equations.  Consider, for example, $\phi^4$ theory in
two dimensions in the broken-symmetry phase.  The equation of motion is
\begin{equation}
  \frac{\partial^2\phi}{\partial t^2} - \nabla^2\phi + \cdots = 0 ,
  \label{eq:22.18}
\end{equation}
which has static solutions, the \emph{kinks}, that interpolate between the
two degenerate vacua.  These solitonic states are the quantum-mechanical
counterparts of the classical solutions.  Similar solitonic excitations, such
as magnetic monopoles and instantons, play important roles in more realistic
theories.

\section{Quantum Field Theory and Gravity}

The final frontier of quantum field theory is the problem of quantizing
gravity itself.  The gravitational interaction is described classically by
Einstein's theory of general relativity.  At the quantum level, the
graviton, the quantum of the gravitational field, is a massless spin-2
particle.  The theory of a massless spin-2 field is highly constrained by
gauge invariance, in a way that is analogous to the constraints on
non-Abelian gauge theories that we studied in Chapter~15.

However, the quantization of gravity runs into a fundamental difficulty: the
theory is not renormalizable by power counting.  The gravitational coupling
constant, Newton's constant $G_N$, has dimensions of inverse mass squared,
\begin{equation}
  G_N \sim \frac{1}{M_{Pl}^2},
  \qquad M_{Pl} \sim 10^{19}\ \mathrm{GeV} ,
  \label{eq:22.19}
\end{equation}
where $M_{Pl}$ is the Planck mass.  By the analysis of Chapter~10, a theory
with a coupling of negative mass dimension is nonrenormalizable: the
divergences cannot be removed by a finite number of counterterms.  Thus, the
quantum theory of gravity, treated as a quantum field theory, has an
infinite number of ultraviolet divergences.

Despite this difficulty, it is still possible to use quantum field theory to
make predictions in the regime where the gravitational interaction is weak.
Because the coupling is so small, the leading-order effects of quantum
gravity are computable: they are corrections to the gravitational potential
of the form
\begin{equation}
  V(r) = -\frac{G_N m_1 m_2}{r}\ \Bigl[ 1 + c\,
    \frac{G_N\hbar}{r^2} + \cdots \Bigr] .
  \label{eq:22.20}
\end{equation}
These corrections are, however, far too small to be measured in any
foreseeable experiment.

The resolution of the problem of quantum gravity requires new physics at the
Planck scale.  One promising approach is string theory, in which the
fundamental objects are not point particles but one-dimensional extended
objects.  In string theory, the graviton appears naturally as one of the
excitations of the string, and the ultraviolet divergences of point-particle
field theory are softened by the extended nature of the string.  At low
energies, string theory reproduces the predictions of quantum field theory,
but at energies of order the Planck mass, its structure is completely
different.

The exploration of quantum gravity remains one of the great frontiers of
theoretical physics.  It is a problem that combines the deepest ideas of
quantum mechanics, general relativity, and quantum field theory, and its
solution will require fundamentally new concepts.  We hope that the tools of
quantum field theory that we have developed in this book will provide the
foundation for the next generation of physicists who will address these
questions.

The study of quantum field theory is far from finished.  The methods
developed in this book have proven to be enormously successful in describing
the behavior of elementary particles, and they continue to be the essential
tools for exploring the frontiers of physics.  We have tried, in this final
chapter, to convey some sense of the excitement and challenge of the current
research frontier.  The problems that remain---the nature of the Higgs
sector, the origin of the fermion masses and mixings, the problem of the
cosmological constant, the quantization of gravity---are among the deepest
questions in all of science.  Quantum field theory, with its remarkable
combination of mathematical elegance and experimental success, will surely
play a central role in addressing them.  We hope that this book has given you
the tools and the inspiration to join in this endeavor.

\problems

\textbf{22.1} \emph{The lattice gauge theory action.}  In this problem we
study the Wilson action for lattice gauge theory.  Show that, in the
continuum limit, the Wilson action for a lattice gauge theory reduces to the
standard Yang-Mills action, and compute the leading corrections.

\textbf{22.2} \emph{Confinement in strong-coupling lattice gauge theory.}
In this problem, we derive the area law for the Wilson loop in the
strong-coupling limit of lattice gauge theory.  Perform the strong-coupling
expansion of the Wilson loop expectation value and show that it decays as
$e^{-\sigma A}$, with a computable string tension.

\textbf{22.3} \emph{Baryon number violation in grand unified theories.}
In this problem we compute the rate of proton decay in the minimal $SU(5)$
grand unified theory.  Show that the decay $p\to e^{+}\pi^0$ is mediated by
the exchange of the heavy $X$ and $Y$ gauge bosons, and compute the lifetime
of the proton in terms of the grand unification scale.
